\documentclass[11pt,a4paper]{article}
\usepackage[utf8]{inputenc}
\usepackage[T1]{fontenc}
\usepackage[french]{babel}
\usepackage{amsmath,amssymb,amsthm}
\usepackage{mathtools}
\usepackage{geometry}
\usepackage{booktabs}
\usepackage{hyperref}
\geometry{margin=2.5cm}

\title{Architecture et formalisation mathematique\\de la tete de regles}
\author{Projet de recherche}
\date{\today}

\begin{document}
\maketitle

\begin{abstract}
Ce document presente de maniere autonome l'architecture de la tete de regles
utilisee dans notre modele hybride. Nous detaillons le pipeline de donnees binaires,
la formalisation logique dure, son approximation differentiable, la couche de sortie
multiclasse, les objectifs d'optimisation, ainsi que la procedure d'extraction de regles
interpretables.y
\end{abstract}

\section{Objectif de la tete de regles}
La tete de regles vise a produire une prediction multiclasse interpretable a partir d'entrees binaires.
L'idee est de representer chaque neurone cache comme une regle candidate (conjonction de litteraux),
puis d'agreger ces regles pour obtenir les logits de classes.

\section{Notations}
Soit:
\begin{itemize}
    \item $\bm{x}\in\mathbb{R}^{d}$ une entree continue;
    \item $B(\cdot)$ l'operateur de binarisation/discretisation;
    \item $\bm{x}^{(b)}=B(\bm{x})\in\{0,1\}^{D_b}$ la representation binaire;
    \item $K$ le nombre de neurones-regles;
    \item $C$ le nombre de classes;
    \item $W_r\in\mathbb{R}^{D_b\times K}$, $\bm{b}_r\in\mathbb{R}^{K}$ les poids de la couche regles;
    \item $W_o\in\mathbb{R}^{K\times C}$, $\bm{b}_o\in\mathbb{R}^{C}$ les poids de la couche de sortie.
\end{itemize}

On regroupe:
\[
W_m=\{W_r,\bm{b}_r,W_o,\bm{b}_o\}.
\]

\section{Architecture}
\subsection{Etape 1: projection binaire}
Les variables continues sont transformees en bits via discretisation (bins) + one-hot:
\[
\bm{x}^{(b)}=B(\bm{x}).
\]
Chaque variable d'origine est remplacee par un sous-vecteur binaire.

\subsection{Etape 2: couche de regles}
La couche cache produit un vecteur d'activations de regles:
\[
\bm{o}=(o_1,\dots,o_K).
\]
Chaque $o_k$ mesure a quel point la regle $k$ est satisfaite.

\subsection{Etape 3: couche OR/additive multiclasse}
Les logits sont obtenus par combinaison lineaire des activations de regles:
\[
\bm{z}=\bm{o}^\top W_o + \bm{b}_o,\qquad
\hat{\bm{p}}=\text{softmax}(\bm{z}).
\]
Prediction finale:
\[
\hat{y}=\arg\max_{c\in\{1,\dots,C\}}\hat{p}_c.
\]

\section{Formalisation mathematique de la couche regles}
\subsection{Version logique dure (ideal theorique)}
Dans l'esprit HyperLogic/DR-Net, pour la regle $k$:
\begin{equation}
o_k
=\mathbf{1}\!\left\{
\sum_{d=1}^{D_b} w_{d,k}x_d^{(b)}-\sum_{d=1}^{D_b}|w_{d,k}|=0
\right\}.
\label{eq:hard_rule}
\end{equation}

Interpretation:
\begin{itemize}
    \item $w_{d,k}>0$: la regle prefere l'activation du bit $x_d^{(b)}$;
    \item $w_{d,k}<0$: la regle prefere la non-activation du bit;
    \item $w_{d,k}\approx 0$: litteral peu informatif pour cette regle.
\end{itemize}

\subsection{Version differentiable (utilisee en apprentissage)}
Comme l'indicatrice de \eqref{eq:hard_rule} n'est pas differentiable, on introduit:
\[
u_k=\sum_{d=1}^{D_b} w_{d,k}x_d^{(b)}-\sum_{d=1}^{D_b}|w_{d,k}|.
\]
Puis une approximation lisse:
\begin{equation}
o_k \approx h(u_k),\qquad
h(u)=\exp\!\left(-\frac{u^2}{\tau}\right),\ \tau>0.
\label{eq:smooth_rule}
\end{equation}

Alternative usuelle en implementation:
\[
o_k=\sigma\!\left(\frac{u_k+b_{r,k}}{\tau}\right),
\]
avec $\sigma$ sigmoide.
Le parametre $\tau$ controle la nettete logique.

\section{Formalisation multiclasse}
Pour chaque classe $c$:
\[
z_c=\sum_{k=1}^{K} u_{k,c}\,o_k+b_{o,c},
\]
ou $u_{k,c}$ est l'element $(k,c)$ de $W_o$.
Chaque classe est donc une combinaison de regles apprises.

\section{Objectif d'optimisation de la tete de regles}
La perte de classification:
\[
\mathcal{L}_{cls}
=\frac{1}{|\mathcal{B}|}\sum_{i\in\mathcal{B}}
\text{CE}\!\left(\hat{\bm{p}}_i,y_i\right).
\]

Regularisation de parcimonie sur la structure regles:
\[
\mathcal{R}_{rule}=\|W_r\|_1.
\]

Objectif total (version tete regles seule):
\[
\mathcal{L}_{rule}
=\mathcal{L}_{cls}+\lambda_{rule}\mathcal{R}_{rule}.
\]

\section{Extraction de regles interpretable}
\subsection{Regle neurone $k$}
On lit la colonne $W_r[:,k]$:
\begin{itemize}
    \item signe du poids $\to$ ON/OFF du litteral;
    \item module du poids $\to$ importance locale du litteral.
\end{itemize}

Une regle textuelle est construite sous la forme:
\[
\text{IF } \ell_{k,1}\wedge \ell_{k,2}\wedge \cdots \wedge \ell_{k,m}
\text{ THEN class } c^\star_k,
\]
avec:
\[
c^\star_k=\arg\max_{c} W_o[k,c].
\]

\subsection{Version compacte vs complete}
\begin{itemize}
    \item \textbf{Compacte}: top-$m$ litteraux selon $|w_{d,k}|$.
    \item \textbf{Complete}: tous les litteraux au-dessus d'un seuil (ou tous, seuil nul).
\end{itemize}

\section{Complexite et compromis}
Augmenter $K$ enrichit le vocabulaire de regles mais peut:
\begin{itemize}
    \item augmenter le risque de redondance;
    \item diminuer la lisibilite globale sans regularisation;
    \item alourdir le cout d'interpretation humaine.
\end{itemize}

Le compromis pratique est pilote par:
\begin{itemize}
    \item $K$ (nombre de regles),
    \item $\lambda_{rule}$ (sparsite),
    \item $\tau$ (rigidite logique).
\end{itemize}

\section{Resume operationnel}
La tete de regles suit la chaine:
\[
\bm{x}\xrightarrow[]{B}\bm{x}^{(b)}
\xrightarrow[]{W_r,\bm{b}_r}\bm{o}
\xrightarrow[]{W_o,\bm{b}_o}\bm{z}
\xrightarrow[]{\text{softmax}}\hat{\bm{p}}
\xrightarrow[]{\arg\max}\hat{y}.
\]

Elle fournit simultanement:
\begin{enumerate}
    \item une prediction multiclasse;
    \item une decomposition interpretable en regles neurones.
\end{enumerate}

\section{Conclusion}
La formalisation precedente montre que la tete de regles est un bloc neuro-symbolique differentiable:
la logique est approchee de maniere lisse pendant l'apprentissage, puis lue symboliquement
au travers des poids pour produire des regles IF-THEN exploitables.

\end{document}
